\chapter{Thiết kế và hiện thực\label{chapter:design}}
\chaptermark{Thiết kế}

Không gì được xây cho đến khi đặc tả được cố định. Nó chốt giao diện lập trình công khai, máy
trạng thái UE, bố cục dữ liệu và các bước trên đó, khoá và kích thước phân mảnh, và mô hình
luồng. Cả hai bản đều bám theo nó. Phép thử khi đó đo ngôn ngữ, không phải một khoảng cách thiết
kế hay kỹ năng. Bản thân đặc tả là một đầu ra của nghiên cứu. Nó là kiểm soát chính chống thiên
lệch của người viết.

Một mô hình luồng được dùng ở cả hai phía: không chia sẻ trạng thái và một luồng mỗi phân mảnh.
Bố cục đó được biết là tự nó đổi độ trễ đuôi~\cite{enberg2019threadpercore}. Không gian định
danh được tách thành các phân mảnh. Mỗi phân mảnh thuộc một luồng xử lý. Chỉ luồng đó đọc hoặc
đổi nó. Vì vậy không khoá nào canh trạng thái ngữ cảnh. Bên gọi gửi một yêu cầu tới hàng đợi bị
chặn của phân mảnh sở hữu. Chúng chờ trên một tay cầm trả lời một lần. Hàng đợi chặn khi đầy.
Điều đó là cố ý. Một quy tắc bỏ sẽ che quá tải. Nghiên cứu này cần quá tải hiện ra như độ trễ.
Trạng thái khả biến dùng chung duy nhất mà bản nào cho phép là một bộ đếm nguyên tử. Không luồng
điều khiển nào phụ thuộc giá trị của nó. Hình~\ref{fig:arch} cho thấy luồng dữ liệu. Có một hàng
đợi nhiều tác nhân sinh, một tác nhân tiêu thụ (MPSC) mỗi phân mảnh. Luồng xử lý của nó rút theo
vào trước ra trước (FIFO).

\begin{figure}[ht]
\centering
\begin{tikzpicture}[
    font=\small,
    prod/.style={draw, rounded corners, minimum width=2.0cm, minimum height=0.66cm, fill=blue!8},
    queue/.style={draw, minimum width=2.5cm, minimum height=0.66cm, fill=gray!12},
    worker/.style={draw, rounded corners, minimum width=2.1cm, minimum height=0.66cm, fill=red!8},
    store/.style={draw, cylinder, shape border rotate=90, aspect=0.25, minimum width=1.5cm, minimum height=0.85cm, fill=green!8},
    lane/.style={draw, dotted, rounded corners, inner sep=7pt},
    >={Stealth[length=2.4mm]}
  ]
  \node[prod] (p1) at (0,1.3)  {Tác nhân 1};
  \node[prod] (p2) at (0,0.0)  {Tác nhân 2};
  \node[prod] (p3) at (0,-1.3) {Tác nhân $m$};
  \begin{scope}[on background layer]
    \node[draw, dashed, rounded corners, fill=blue!3, fit=(p1)(p2)(p3), inner sep=6pt] (pg) {};
  \end{scope}
  \node[font=\scriptsize] at (0,-2.15) {bộ sinh tải};

  \node[queue] (q1) at (5.2,1.3)   {Hàng đợi MPSC};
  \node[worker](w1) at (8.9,1.3)   {Worker 1};
  \node[store] (s1) at (12.0,1.3)  {Phân mảnh 1};
  \node[queue] (qn) at (5.2,-1.3)  {Hàng đợi MPSC};
  \node[worker](wn) at (8.9,-1.3)  {Worker $n$};
  \node[store] (sn) at (12.0,-1.3) {Phân mảnh $n$};
  \begin{scope}[on background layer]
    \node[lane, fit=(q1)(s1)] (lane1) {};
    \node[lane, fit=(qn)(sn)] (lanen) {};
  \end{scope}
  \node[font=\scriptsize] at (8.35,2.35) {luồng sở hữu phân mảnh: không trạng thái khả biến dùng chung, không khoá};

  \node[font=\scriptsize, align=center, fill=white, inner sep=2pt] at (2.85,0) {định tuyến theo\\ hash(\texttt{ue\_id})};
  \draw[->] (pg.east) -- (q1.west) node[midway, above, sloped, font=\scriptsize]{gửi};
  \draw[->] (pg.east) -- (qn.west) node[midway, below, sloped, font=\scriptsize]{gửi};

  \draw[->]  (q1) -- node[above, font=\scriptsize]{FIFO} (w1);
  \draw[->]  (qn) -- node[above, font=\scriptsize]{FIFO} (wn);
  \draw[<->] (w1) -- node[above, font=\scriptsize]{sở hữu} (s1);
  \draw[<->] (wn) -- node[above, font=\scriptsize]{sở hữu} (sn);

  \draw[->, dashed, rounded corners=4pt]
      (wn.south) -- ++(0,-0.75) -- (1.4,-2.9) -- ($(pg.south east)+(-0.4,0)$);
  \node[font=\scriptsize] at (5.9,-3.05) {hoàn tất một lần (kết quả $\rightarrow$ tác nhân đang chờ)};
\end{tikzpicture}
\caption[Kiến trúc không chia sẻ, một luồng mỗi phân mảnh]{Kiến trúc không chia sẻ, một luồng
mỗi phân mảnh. Các tác nhân sinh gửi vào một hàng đợi MPSC bị chặn mỗi phân mảnh, và luồng xử
lý của phân mảnh đó rút hàng.\label{fig:arch}}
\end{figure}

Tính công bằng được giữ dưới thiết kế cũng như trong thiết kế. Cả hai bản build liên kết cùng
\texttt{jemalloc}. Cả hai giữ một bản build bộ cấp phát hệ thống làm đối chứng. Cả hai biên dịch
ở cùng mức tối ưu. Cả hai ghi độ trễ qua một bộ histogram dải động cao (HDR) được lập cùng một
cách~\cite{tene_hdrhistogram}. Băm được chốt ở cả hai phía. SplitMix64 gán phân mảnh. FNV-1a,
hàm băm Fowler--Noll--Vo, được dùng bên trong bảng. Hàm băm mặc định của C++ trên một số nguyên
thường chỉ là chính số đó. Của Rust mạnh là cố ý. Để nguyên như chúng đến, hai thư viện chuẩn
sẽ là phép đo.

Nơi các ngôn ngữ không thể đáp một nhu cầu cùng một cách, khoảng cách được viết vào nhật ký
thành ngữ bắt buộc rồi được đo. Hai trong số các khoảng cách đó mang kết quả quá tải của
Mục~\ref{sec:overload}. Một là hàng đợi bị chặn. Rust lấy nó từ \texttt{sync\_channel}. C++23
vẫn chưa có nguyên thủy đó, nên C++ xây một cái từ một mutex và hai biến điều kiện. Cái kia là
tay cầm trả lời, \texttt{std::promise} đối với một kênh \texttt{mpsc} mới mỗi yêu cầu. Một
khoảng cách thứ ba mang kết quả thao tác lõi của Mục~\ref{sec:perfresults}: container mỗi phân
mảnh là bảng băm chuẩn của chính ngôn ngữ đó. Một khoảng cách thứ tư là mã hoá yêu cầu. C++
không có kiểu liệt kê mang payload theo thành ngữ này, nên một yêu cầu thành một struct có nhãn
giữ mỗi trường mà thao tác nào cũng có thể cần, $\CppRequestBytes$~byte đối với
$\RustRequestBytes$ của Rust. Kết quả ròng là một lõi Rust dài hơn $\NlocRise\%$ về dòng mã. Độ
phức tạp cyclomatic của nó thấp hơn $\CcnTotalDrop\%$~\cite{lizard}. Thư viện Rust không chứa
\texttt{unsafe} nào cả.

Tương đương được thử chứ không được giả định. Các vector trung lập ngôn ngữ chạy ở cả hai bản.
Nhật ký của chúng được diff theo từng thao tác. Một diff chéo không phân biệt ``cả hai đúng''
với ``cả hai sai cùng một cách''. Một bản thứ ba bằng Python đóng vai tiên tri độc lập. Nó phải
tái tạo tất cả $\ConfHandOps$ kỳ vọng suy bằng tay trước khi được sinh kỳ vọng của riêng mình.
Một chiến dịch $\ConfGenOps$ thao tác được sinh không tìm thấy lệch. Phía C++ chạy dưới
AddressSanitizer, ThreadSanitizer và UndefinedBehaviorSanitizer. Phía Rust chạy dưới Miri. Một
kết quả bộ dò lỗi sạch bị chặn bởi mã thực sự đã chạy.
